\section{Results}
\label{sec:results}

\subsection{Behavioral Evaluation}
\label{sec:results_behavioral}

\Tabref{tab:behavioral} presents the behavioral results for \gptmodel.
The model achieves \textbf{100\% accuracy} on every task, with no difference between vigesimal and decimal numbers.
All six counting sequences---three crossing vigesimal boundaries and three decimal controls---produce perfectly correct French number words with correct transitions.

\begin{table}[t]
    \centering
    \caption{Behavioral accuracy of \gptmodel on French number tasks. The model achieves perfect accuracy across all conditions, including vigesimal numbers (70--99) that require implicit base-20 arithmetic.}
    \begin{tabular}{@{}lccc@{}}
        \toprule
        \textbf{Task} & \textbf{Decimal (0--69)} & \textbf{Vigesimal (70--99)} & \textbf{Overall} \\
        \midrule
        French $\to$ Number & 100\% & 100\% & 100\% \\
        Number $\to$ French & 100\% & 100\% & 100\% \\
        Belgian French $\to$ Number & --- & 100\% & 100\% \\
        Arithmetic & 100\% & 100\% & 100\% \\
        Counting sequences & 100\% & 100\% & 100\% \\
        \bottomrule
    \end{tabular}
    \label{tab:behavioral}
\end{table}

This result refutes the hypothesis that vigesimal numbers cause behavioral difficulties for frontier models.
\gptmodel has fully internalized the French counting system from its training data, including all vigesimal irregularities.

\subsection{Tokenization Overhead}
\label{sec:results_tokens}

Vigesimal French numbers require substantially more tokens than their decimal counterparts (\tabref{tab:tokens}).
The mean token count for vigesimal words is 7.3, compared to 5.2 for decimal French words---a 40\% increase.
At the extreme, \emph{quatre-vingt-dix-sept} (97) requires 10 tokens, compared to just 3 for the digit string ``97.''
Digit representations are constant at 3.0 tokens regardless of vigesimal status.
This increased token count means the model must propagate numeric information through more transformer steps for vigesimal numbers.

\begin{table}[t]
    \centering
    \caption{Mean token counts for \mistral's tokenizer across number forms and ranges. Vigesimal French words require 40\% more tokens than decimal French words.}
    \begin{tabular}{@{}lcc@{}}
        \toprule
        \textbf{Number Form} & \textbf{Decimal (20--69)} & \textbf{Vigesimal (70--99)} \\
        \midrule
        French word & 5.2 & \textbf{7.3} (+40\%) \\
        Digit string & 3.0 & 3.0 \\
        \bottomrule
    \end{tabular}
    \label{tab:tokens}
\end{table}

\subsection{Ordinal Structure of Representations}
\label{sec:results_ordinal}

We measure how well each form's representational distances preserve numeric ordering using Spearman correlation between the RDM and the ideal numeric distance matrix (\tabref{tab:ordinal}).
Belgian French representations achieve the highest ordinal correlation ($\rho = 0.591$, $p < 10^{-300}$), surpassing both standard French ($\rho = 0.520$) and, notably, digit representations ($\rho = 0.485$).
The regular decimal system of Belgian French thus produces internal representations that are more numerically organized than either the vigesimal system or raw digits.

\begin{table}[t]
    \centering
    \caption{Ordinal correlation (Spearman $\rho$) between representational distances and numeric distances at the final layer of \mistral. Belgian French, with its regular decimal system, achieves the best ordinal structure. All $p < 10^{-280}$.}
    \begin{tabular}{@{}lc@{}}
        \toprule
        \textbf{Form} & \textbf{Spearman $\rho$} \\
        \midrule
        French & 0.520 \\
        Digit & 0.485 \\
        Belgian French & \textbf{0.591} \\
        \bottomrule
    \end{tabular}
    \label{tab:ordinal}
\end{table}

\subsection{Linear Probe Accuracy}
\label{sec:results_probe}

\Tabref{tab:probe} shows Ridge regression probe results at the final layer.
French number words produce the most linearly decodable representations ($R^2 = 0.979$, $\MAE = 2.72$), followed closely by Belgian French ($R^2 = 0.977$, $\MAE = 2.79$), with digit strings substantially behind ($R^2 = 0.943$, $\MAE = 4.18$).
Within each form, vigesimal numbers have higher $\MAE$ than decimal numbers, confirming that the vigesimal range is harder to decode linearly.

\begin{table}[t]
    \centering
    \caption{Linear probe results (Ridge regression, 5-fold CV) predicting numeric value from \mistral final-layer hidden states. French words are more linearly decodable than digit strings. Best results in \textbf{bold}.}
    \begin{tabular}{@{}lcccc@{}}
        \toprule
        \textbf{Form} & \textbf{MAE (overall)} & \textbf{$R^2$} & \textbf{Vigesimal MAE} & \textbf{Decimal MAE} \\
        \midrule
        French & \textbf{2.72 $\pm$ 0.65} & \textbf{0.979} & 2.37 & 1.82 \\
        Belgian French & 2.79 $\pm$ 0.53 & 0.977 & 2.54 & 1.77 \\
        Digit & 4.18 $\pm$ 1.19 & 0.943 & 4.37 & 3.49 \\
        \bottomrule
    \end{tabular}
    \label{tab:probe}
\end{table}

The finding that contextualized French words are more decodable than digits is counterintuitive.
We attribute this to the prompt template ``\emph{Le nombre est [X]},'' which activates a numeric semantic frame, producing richer representations than raw digit strings.

\subsection{Vigesimal Clustering by Linguistic Base}
\label{sec:results_clustering}

The central question of our representational analysis is whether vigesimal numbers cluster by linguistic base or numeric value.
\Tabref{tab:nn} presents the nearest-neighbor analysis for numbers 70--99 in French representation space.

\begin{table}[t]
    \centering
    \caption{Nearest-neighbor analysis for vigesimal French numbers (70--99). The majority of nearest neighbors share a linguistic base rather than being numerically adjacent. Examples show numbers whose representational nearest neighbor is numerically distant but linguistically similar.}
    \begin{tabular}{@{}lccl@{}}
        \toprule
        \textbf{Metric} & \textbf{Vigesimal (70--99)} & \textbf{Decimal (20--69)} \\
        \midrule
        NN shares linguistic base & \textbf{70\%} & --- \\
        NN within $\pm 2$ numerically & 53.3\% & 34\% \\
        Mean NN error & 4.3 $\pm$ 4.3 & 2.7 $\pm$ 1.7 \\
        \midrule
        \multicolumn{3}{@{}l}{\textbf{Notable clustering examples:}} \\
        \midrule
        \multicolumn{3}{@{}l}{80 (\emph{quatre-vingts}) $\to$ NN: 95 (\emph{quatre-vingt-quinze}), $\Delta = 15$} \\
        \multicolumn{3}{@{}l}{83 (\emph{quatre-vingt-trois}) $\to$ NN: 93 (\emph{quatre-vingt-treize}), $\Delta = 10$} \\
        \multicolumn{3}{@{}l}{73 (\emph{soixante-treize}) $\to$ NN: 63 (\emph{soixante-trois}), $\Delta = 10$} \\
        \multicolumn{3}{@{}l}{91 (\emph{quatre-vingt-onze}) $\to$ NN: 82 (\emph{quatre-vingt-deux}), $\Delta = 9$} \\
        \bottomrule
    \end{tabular}
    \label{tab:nn}
\end{table}

\textbf{70\% of vigesimal numbers' nearest neighbors share a linguistic base}---\eg both 83 (\emph{quatre-vingt-trois}) and 93 (\emph{quatre-vingt-treize}) begin with \quatrevingt and are representationally close despite being 10 apart numerically.
Similarly, 73 (\emph{soixante-treize}) and 63 (\emph{soixante-trois}) cluster together because they share the \emph{soixante-} prefix.
The pattern is bimodal: many vigesimal numbers do follow numeric ordering (NN distance = 1), but a substantial fraction have nearest neighbors driven by shared prefixes.

A Mann-Whitney U test comparing vigesimal vs.\ decimal NN distances yields $U = 708.0$, $p = 0.667$ (not significant), indicating that vigesimal numbers are not \emph{uniformly} worse---rather, the effect is concentrated in specific numbers where linguistic similarity overrides numeric proximity.

\subsection{Cross-Form Similarity}
\label{sec:results_crossform}

\Tabref{tab:crossform} shows mean cosine similarity between representations of the same number across different forms.
French and Belgian French representations are nearly identical ($0.995$ cosine similarity), suggesting the model maps both surface forms to a shared underlying representation.
Both French forms are moderately aligned with digit representations ($0.907$).

\begin{table}[t]
    \centering
    \caption{Mean cosine similarity between representations of the same number across different forms in \mistral's final layer. French and Belgian French map to nearly identical internal representations despite different surface forms.}
    \begin{tabular}{@{}lc@{}}
        \toprule
        \textbf{Comparison} & \textbf{Cosine Similarity} \\
        \midrule
        French--Digit & 0.907 \\
        Belgian--Digit & 0.907 \\
        French--Belgian & \textbf{0.995} \\
        \bottomrule
    \end{tabular}
    \label{tab:crossform}
\end{table}

\subsection{Representational Similarity Analysis}
\label{sec:results_rsa}

\Tabref{tab:rsa} reports Spearman correlations between RDMs.
Belgian French's RDM is most correlated with the ideal numeric RDM ($\rho = 0.591$), confirming its superior ordinal structure.
The French--Belgian RDM correlation ($\rho = 0.890$) is high but not as extreme as the per-number cosine similarity ($0.995$), indicating that while individual number representations are nearly identical, the \emph{pairwise distance structure} differs---consistent with the linguistic clustering effect in the vigesimal range.

\begin{table}[t]
    \centering
    \caption{Representational Similarity Analysis (RSA): Spearman $\rho$ between representational dissimilarity matrices (RDMs). Belgian French has the strongest alignment with ideal numeric ordering.}
    \begin{tabular}{@{}lc@{}}
        \toprule
        \textbf{Comparison} & \textbf{Spearman $\rho$} \\
        \midrule
        French vs.\ Ideal Numeric & 0.520 \\
        Digit vs.\ Ideal Numeric & 0.485 \\
        Belgian vs.\ Ideal Numeric & \textbf{0.591} \\
        \cmidrule(lr){1-2}
        French vs.\ Digit & 0.656 \\
        Belgian vs.\ Digit & 0.661 \\
        French vs.\ Belgian & \textbf{0.890} \\
        \bottomrule
    \end{tabular}
    \label{tab:rsa}
\end{table}

\subsection{Layer Progression}
\label{sec:results_layers}

\Tabref{tab:layers} tracks how number representations develop across \mistral's layers.
Probe accuracy improves throughout the network, with the best performance at the final layer ($\MAE = 15.0$).
The ordinal correlation peaks at early-to-middle layers ($\rho = 0.526$ at layer 8) before dipping in later layers and partially recovering.
This pattern is consistent with the model processing linguistic features in early layers and reconstructing task-relevant numeric information in later layers.

\begin{table}[t]
    \centering
    \caption{Layer-wise probe accuracy and ordinal correlation for French number representations in \mistral. Numeric information develops progressively across layers.}
    \begin{tabular}{@{}lcc@{}}
        \toprule
        \textbf{Layer} & \textbf{Probe MAE} & \textbf{Spearman $\rho$} \\
        \midrule
        0 (embedding) & 31.0 & --- \\
        8 & 18.8 & \textbf{0.526} \\
        16 (middle) & 16.4 & 0.525 \\
        24 & 18.4 & 0.490 \\
        32 (final) & \textbf{15.0} & 0.520 \\
        \bottomrule
    \end{tabular}
    \label{tab:layers}
\end{table}
